% ============================================================================
% SPANISCH LANGENTWURF - DS 1 BUCHARBEIT PARALLEL
% ============================================================================
% Basiert auf: /templates/spanisch/langentwurf-spanisch.tex
% Fachfarbe: #c41e3a (Karminrot)
%
% KONTEXT: Vorbereitung Unterrichtsbeobachtung Februar 2026
% KURS: Jahrgangsübergreifend Spanisch Spätbeginner (Kl. 10, 11, 12)
%
% DIFFERENZIERUNG:
%   [LJ1] = Gruppe A (Lernjahr 1, A2-Niveau) - 7 SuS
%   [LJ2/LJ3] = Gruppe B (Lernjahr 2-3, B1-Niveau) - 3 SuS
%
% PFLICHT: Tier 1 (vollständig) für alle Doppelstunden
% ============================================================================

\documentclass[11pt,a4paper]{article}

% ============================================================================
% PACKAGES
% ============================================================================
\usepackage[utf8]{inputenc}
\usepackage[T1]{fontenc}
\usepackage[ngerman]{babel}
\usepackage{geometry}
\usepackage{fancyhdr}
\usepackage{lastpage}
\usepackage{xcolor}
\usepackage{tcolorbox}
\usepackage{enumitem}
\usepackage{tabularx}
\usepackage{longtable}
\usepackage{array}
\usepackage{booktabs}
\usepackage{hyperref}
\usepackage{ifthen}
\usepackage{amssymb}

% ============================================================================
% KONFIGURATION
% ============================================================================

% --- TIER (immer 1 für Spanisch) ---
\newcommand{\plantier}{1}

% --- FACH ---
\newcommand{\fachid}{spanisch}
\newcommand{\fach}{Spanisch}
\newcommand{\fachfarbe}{c41e3a}

% ============================================================================
% VARIABLEN
% ============================================================================

% --- Metadaten ---
\newcommand{\jahrgangsstufe}{10--12}
\newcommand{\schulart}{Gesamtschule}
\newcommand{\stundenthema}{Bucharbeit parallel (A: S.30--31 / B: S.102--103)}
\newcommand{\reihenthema}{Beobachtungsvorbereitung Feb 2026}
\newcommand{\stundennummer}{1}
\newcommand{\reihenstunden}{10}
\newcommand{\unterrichtsdatum}{05.01.2026}
\newcommand{\unterrichtszeit}{[laut Stundenplan]}
\newcommand{\raum}{[Kursraum]}

% --- Lehrkraft ---
\newcommand{\lehrkraft}{[N.N.]}
\newcommand{\ausbildungsstand}{Lehrkraft}

% --- Lerngruppe ---
\newcommand{\lerngruppengroesse}{10}
\newcommand{\maennlich}{1}
\newcommand{\weiblich}{9}
\newcommand{\divers}{0}

% --- Kompetenzen/Niveau ---
\newcommand{\gerniveau}{A2/B1}
\newcommand{\lehrwerk}{A\_tope.com}
\newcommand{\lehrwerkseite}{S. 30--31 (A) / 102--103 (B)}

% ============================================================================
% LAYOUT
% ============================================================================
\geometry{left=2.5cm, right=2.5cm, top=2.5cm, bottom=2.5cm}
\definecolor{fach}{HTML}{\fachfarbe}
\definecolor{gray}{HTML}{666666}
\definecolor{lightgray}{HTML}{f5f5f5}
\definecolor{successgreen}{HTML}{2a7a4b}
\definecolor{warningorange}{HTML}{b35c00}

\tcbuselibrary{skins,breakable}

% Farbige Boxen
\newtcolorbox{infobox}[1][]{
    colback=lightgray,
    colframe=fach,
    fonttitle=\bfseries,
    title=#1,
    breakable
}

\newtcolorbox{scorebox}{
    colback=white,
    colframe=fach,
    fonttitle=\bfseries,
    title=Didaktik-Score,
    breakable
}

% Kopf-/Fußzeile
\pagestyle{fancy}
\fancyhf{}
\fancyhead[L]{\small\textcolor{gray}{\fach{} -- Jg. \jahrgangsstufe{} -- \stundenthema}}
\fancyhead[R]{\small\textcolor{gray}{Langentwurf Tier 1}}
\fancyfoot[C]{\small\thepage{} / \pageref{LastPage}}
\renewcommand{\headrulewidth}{0.4pt}
\renewcommand{\footrulewidth}{0pt}

% ============================================================================
% DOKUMENT
% ============================================================================
\begin{document}

% ============================================================================
% DECKBLATT
% ============================================================================
\begin{titlepage}
    \centering
    \vspace*{2cm}

    {\large\textcolor{gray}{\schulart}}\\[2cm]

    {\Huge\textcolor{fach}{\textbf{Langentwurf}}}\\[0.5cm]
    {\large Tier 1 -- Vollständig}\\[2cm]

    {\LARGE\textbf{\stundenthema}}\\[0.5cm]
    {\large im Rahmen der Unterrichtsreihe}\\[0.3cm]
    {\Large\textit{\reihenthema}}\\[2cm]

    \begin{tabular}{rl}
        \textbf{Fach:} & \fach{} (\gerniveau) \\[0.3cm]
        \textbf{Klasse:} & \jahrgangsstufe{} (jahrgangsübergreifend) \\[0.3cm]
        \textbf{Lehrwerk:} & \lehrwerk, \lehrwerkseite \\[0.3cm]
        \textbf{Datum:} & \unterrichtsdatum \\[0.3cm]
        \textbf{Zeit:} & \underline{\hspace{3cm}} (Doppelstunde, 90 Min) \\[0.3cm]
        \textbf{Raum:} & \underline{\hspace{3cm}} \\[0.3cm]
    \end{tabular}

    \vfill

    {\large Lehrkraft: \lehrkraft}\\[0.3cm]
    {\normalsize\textcolor{gray}{\ausbildungsstand}}\\[1cm]

\end{titlepage}

% ============================================================================
% INHALTSVERZEICHNIS
% ============================================================================
\tableofcontents
\newpage

% ============================================================================
% 1. BEDINGUNGSANALYSE
% ============================================================================
\section{Bedingungsanalyse}

\subsection{Lerngruppe}
\begin{tabular}{ll}
    Kursgröße: & \lerngruppengroesse{} SuS (\maennlich{} m / \weiblich{} w / \divers{} d) \\
    Jahrgangsstufe: & \jahrgangsstufe{} (jahrgangsübergreifend) \\
    Sprachniveau: & \gerniveau{} (GER) -- heterogen \\
    Fremdsprache: & Spanisch als spätbeginnende Fremdsprache \\
    Besonderheit: & Zwei Niveaugruppen arbeiten parallel \\
\end{tabular}

\subsection{Gruppenzusammensetzung}

\subsubsection{Gruppe A -- Lernjahr 1 (A2-Niveau): 7 SuS}

\begin{tabular}{|l|c|p{3.5cm}|p{3.5cm}|p{3.5cm}|}
\hline
\textbf{Kürzel} & \textbf{Niveau} & \textbf{Stärken} & \textbf{Schwächen} & \textbf{Strategie} \\
\hline
E.H. (Kl.10) & $\star\star\star\star$ & Franz. Hintergrund, 15 NP, immer aktiv, analytisch & -- & Peer-Tutor, anspruchsvolle Zusatzaufgaben \\
\hline
L.N. (Kl.10) & $\star\star\star$ & Proaktiv, versteht schnell, gute Aussprache & Etwas faul, leicht abgelenkt & In Bewegung halten, Verantwortung geben \\
\hline
C.P. (Kl.10) & $\star\star$ & Gute Aussprache, 12 NP, zuverlässig & Braucht Zeit zum Verstehen, ruhig & Klare Anweisungen, Wartezeit gewähren \\
\hline
V.S. (Kl.10) & $\star\star$ & Konstant, fleißig, ruhig & Weniger Selbstvertrauen & Kleine Erfolgserlebnisse, gezieltes Lob \\
\hline
L.K. (Kl.11) & $\star\star$ & Versteht gut, 10 NP, catch-up aus Kl.11 & Ablenkbar, weniger fleißig & Klare Struktur, Erwartungen kommunizieren \\
\hline
H.P. (Kl.10) & $\star$ & Bemüht, 10 NP & 1--2 Monate später angefangen, Grundlagenlücken & Scaffolding, Partnering mit E.H. \\
\hline
A.A. (Kl.10) & $\star$ & Fleißig, lernt viel, gute Testergebnisse & Unsicher, schaut zu E.H. & Ermutigung, manchmal getrennt von E.H. \\
\hline
A.S. (Kl.10) & $\star$ & Macht Fortschritte, 9 NP, einziger Junge & Unsicher, Aussagen klingen wie Fragen & Positive Verstärkung, klare Strukturen \\
\hline
\end{tabular}

\subsubsection{Gruppe B -- Lernjahr 2--3 (B1-Niveau): 3 SuS}

\begin{tabular}{|l|c|p{5cm}|p{5cm}|}
\hline
\textbf{Kürzel} & \textbf{Klasse} & \textbf{Besonderheiten} & \textbf{Strategie} \\
\hline
SuS 1 & 12 & LRS-Nachteilsausgleich & Zeitzugabe, mündliche Alternativen \\
\hline
SuS 2 & 11 & Stärker, gute Mitarbeit & Kann Gruppe B leiten \\
\hline
SuS 3 & 12 & Solide, zuverlässig & Eigenständige Arbeit möglich \\
\hline
\end{tabular}

\vspace{0.5em}
\textbf{Sitzordnung:} Gruppe B sitzt hinten im Raum (3 Tische), Gruppe A vorne und mittig verteilt. Dies ermöglicht parallele Arbeit mit minimaler gegenseitiger Störung.

\subsection{Differenzierung (intern)}

\textbf{WICHTIG:} Keine Sternchen auf Schülermaterial! Differenzierung über Gruppenbezeichnung.

\begin{tabular}{|l|p{5cm}|p{5cm}|}
\hline
& \textbf{[LJ1] Gruppe A} & \textbf{[LJ2/LJ3] Gruppe B} \\
\hline
\textbf{Niveau} & A2 (Elementare Sprachverwendung) & B1 (Selbstständige Sprachverwendung) \\
\hline
\textbf{Inhalt} & S. 30--31: ser/estar/hay, vuestro, Stadtviertel & S. 102--103: Schulsystem, Mediación, Konjunktionen \\
\hline
\textbf{Produktion} & Gelenkte Dialoge, Blogeintrag mit Hilfen & Freie Erklärungen, Vergleiche, Mediación \\
\hline
\textbf{Hilfen} & Wortschatzlisten, Beispielsätze, Partnerhilfe & Nur Impulse, eigenständige Recherche \\
\hline
\end{tabular}

\subsection{Besonderheiten}
\begin{itemize}
    \item \textbf{LRS:} 1 SuS in Gruppe B (Nachteilsausgleich: Zeitzugabe, keine Bewertung Rechtschreibung)
    \item \textbf{Späteinsteiger:} H.P. hat Grundlagenlücken durch späteren Kurseintritt
    \item \textbf{Peer-Tutoring:} E.H. kann schwächere SuS unterstützen (aber nicht überstrapazieren)
    \item \textbf{Geschlechterverteilung:} A.S. ist einziger männlicher SuS -- auf Integration achten
\end{itemize}

% ============================================================================
% 2. SACHANALYSE
% ============================================================================
\section{Sachanalyse}

\subsection{Sprachliche Strukturen -- Gruppe A}

\subsubsection{Grammatik: ser/estar/hay}
\begin{itemize}
    \item \textbf{ser:} Identität, Herkunft, Eigenschaften (permanente Charakteristika)
    \begin{itemize}
        \item \textit{Mi barrio es grande. El parque es bonito.}
    \end{itemize}
    \item \textbf{estar:} Ort, Zustand (temporär oder lokativ)
    \begin{itemize}
        \item \textit{El cine está en la calle Mayor. Mi casa está cerca.}
    \end{itemize}
    \item \textbf{hay:} Existenz, unpersönlich (Singular = Plural)
    \begin{itemize}
        \item \textit{Hay un parque. Hay muchas tiendas.}
    \end{itemize}
\end{itemize}

\subsubsection{Grammatik: Possessivpronomen ``vuestro''}
\begin{itemize}
    \item Ergänzung zu \textit{mi, tu, su, nuestro}
    \item Kongruenz mit Bezugswort: \textit{vuestro barrio, vuestra casa, vuestros padres, vuestras amigas}
    \item In Spanien gebräuchlich (in Lateinamerika: \textit{su})
\end{itemize}

\subsubsection{Wortschatz: Stadtviertel und Orte}
\begin{itemize}
    \item Orte: \textit{el parque, la piscina, el cine, el teatro, la biblioteca, el supermercado, la iglesia, el hospital}
    \item Beschreibungen: \textit{grande, pequeño, tranquilo, ruidoso, moderno, antiguo, bonito, feo}
    \item Lage: \textit{cerca (de), lejos (de), al lado de, enfrente de, entre}
\end{itemize}

\subsection{Sprachliche Strukturen -- Gruppe B}

\subsubsection{Thema: Schulsystem Spanien vs. Deutschland}
\begin{itemize}
    \item Spanisches System: Educación Primaria (6--12), ESO (12--16), Bachillerato (16--18), FP (Formación Profesional)
    \item Vergleich mit deutschem System: Grundschule, Sekundarstufe I (Gesamt-/Realschule/Gymnasium), Sekundarstufe II
    \item Unterschiede: Ganztagsschule, Notensystem (0--10 vs. 1--6/15 NP), Abschlussprüfungen
\end{itemize}

\subsubsection{Grammatik: Kausale/konzessive Konjunktionen}
\begin{itemize}
    \item \textbf{ya que} (da, weil -- formeller): \textit{Ya que el sistema es diferente, hay que explicarlo.}
    \item \textbf{porque} (weil -- neutral): \textit{Estudio español porque me gusta.}
    \item \textbf{mientras (que)} (während): \textit{Mientras en Alemania hay varios tipos de escuela, en España...}
    \item \textbf{además} (außerdem): Zur Erweiterung von Argumenten
\end{itemize}

\subsubsection{Mediación (Sprachmittlung)}
\begin{itemize}
    \item Kompetenz: Inhalte sinngemäß in andere Sprache übertragen (nicht wörtlich übersetzen)
    \item Aufg. 7: FP-Flyer auf Spanisch erklären (deutsche Informationen spanisch wiedergeben)
    \item Strategien: Umschreibungen, Vereinfachungen, kulturelle Anpassungen
\end{itemize}

\subsection{Kontrastive Betrachtung (Deutsch $\leftrightarrow$ Spanisch)}

\begin{itemize}
    \item \textbf{Positive Transfers:}
    \begin{itemize}
        \item Lateinische/romanische Lehnwörter im Deutschen erleichtern Wortschatz
        \item Ähnliche Satzstellung in einfachen Aussagesätzen
    \end{itemize}
    \item \textbf{Interferenzen:}
    \begin{itemize}
        \item Ser/estar-Unterscheidung existiert im Deutschen nicht (``sein'' für beides)
        \item Hay als unpersönliche Form schwer zu verinnerlichen (Deutsche suchen Subjekt)
        \item Possessivpronomen: Kongruenz im Spanischen (vuestro/-a/-os/-as) vs. Deutsch (euer)
    \end{itemize}
    \item \textbf{Falsche Freunde:}
    \begin{itemize}
        \item \textit{actual} $\neq$ aktuell (= \textit{tatsächlich})
        \item \textit{mayor} $\neq$ Major (= \textit{größer, älter})
    \end{itemize}
\end{itemize}

\subsection{Kulturelle Aspekte}

\begin{itemize}
    \item \textbf{Stadtviertel in Spanien:} Starke Identifikation mit dem \textit{barrio}, Nachbarschaft als Gemeinschaft
    \item \textbf{Bildungssystem:} Zentralisiert, aber regionale Unterschiede (Comunidades Autónomas)
    \item \textbf{Formación Profesional:} Aufwertung der Berufsausbildung in Spanien seit 2010er Jahren
\end{itemize}

% ============================================================================
% 3. DIDAKTISCHE ANALYSE
% ============================================================================
\section{Didaktische Analyse}

\subsection{Stellung in der Sequenz}
Dies ist die \textbf{\stundennummer. Doppelstunde} der Sequenz ``\reihenthema'' (\reihenstunden{} Doppelstunden gesamt).

\textbf{Sequenzübersicht:}
\begin{enumerate}
    \item DS 1--2: Bucharbeit parallel (Gruppe A: S. 30--34, Gruppe B: S. 102--107)
    \item DS 3--4: Bucharbeit + Assessment-Vorbereitung
    \item DS 5: Mündliches Assessment (sonstige Note)
    \item DS 6--8: Brückenthema + Konvergenz (beide Gruppen zusammen)
    \item DS 9: Puffer
    \item DS 10: \textbf{Beobachtungsstunde} (03.02.2026)
\end{enumerate}

\subsection{Kompetenzziele (Rahmenplan MV Spanisch)}

\begin{itemize}
    \item \textbf{Kommunikative Kompetenz:}
    \begin{itemize}
        \item [A] Sprechen: Dialogisches Sprechen (Partnerarbeit Madrid-Dialog)
        \item [A] Schreiben: Zusammenhängender Text (Blogeintrag Stadtviertel)
        \item [B] Sprechen: Monologisches Sprechen (Erklärung Schulsystem)
        \item [B] Sprachmittlung: Mediación (FP-Flyer)
    \end{itemize}
    \item \textbf{Sprachliche Kompetenz:}
    \begin{itemize}
        \item [A] Grammatik: ser/estar/hay, Possessivpronomen vuestro
        \item [B] Grammatik: Kausale Konjunktionen (ya que, porque, mientras)
    \end{itemize}
    \item \textbf{Interkulturelle Kompetenz:}
    \begin{itemize}
        \item [A] Stadtviertel in spanischsprachigen Ländern
        \item [B] Vergleich Bildungssysteme Spanien/Deutschland
    \end{itemize}
    \item \textbf{Methodische Kompetenz:}
    \begin{itemize}
        \item Selbstständiges Arbeiten mit dem Lehrwerk
        \item Partnerarbeit mit verteilten Rollen (A mit Partner B, S. 136)
    \end{itemize}
\end{itemize}

\subsection{Legitimation}
\textbf{Rahmenplan Spanisch MV (Sek I/II):}
\begin{itemize}
    \item \textbf{Kompetenzbereich:} Funktionale kommunikative Kompetenz
    \item \textbf{Standard A2:} ``Die SuS können in einfachen, routinemäßigen Situationen kommunizieren, in denen es um einen direkten Austausch von Informationen über vertraute und geläufige Dinge geht.''
    \item \textbf{Standard B1:} ``Die SuS können die Hauptpunkte verstehen, wenn klare Standardsprache verwendet wird und wenn es um vertraute Dinge aus Arbeit, Schule, Freizeit usw. geht.''
\end{itemize}

\subsection{Didaktische Reduktion}

\begin{itemize}
    \item \textbf{Gruppe A:} Fokus auf produktive Anwendung von ser/estar/hay (nicht alle Sonderfälle)
    \item \textbf{Gruppe B:} Schulsystem vereinfacht dargestellt, keine historische Entwicklung
    \item \textbf{Beide:} Konzentration auf kommunikative Anwendung, nicht grammatische Metasprache
\end{itemize}

% ============================================================================
% 4. STUNDENZIELE
% ============================================================================
\section{Stundenziele}

\subsection{Grobziel -- Gruppe A [LJ1]}
Die SuS erweitern ihre \textbf{kommunikative Kompetenz} im Bereich Sprechen und Schreiben, indem sie ihr Stadtviertel mithilfe von \textit{ser/estar/hay} beschreiben und einen kurzen Blogeintrag verfassen.

\subsection{Grobziel -- Gruppe B [LJ2/LJ3]}
Die SuS erweitern ihre \textbf{kommunikative Kompetenz} im Bereich monologisches Sprechen und Sprachmittlung, indem sie das spanische Schulsystem erklären, mit dem deutschen vergleichen und einen FP-Flyer sinngemäß ins Spanische übertragen.

\subsection{Feinziele (operationalisiert)}

\textbf{Gruppe A:}

\begin{tabular}{|c|p{7cm}|c|c|}
\hline
\textbf{Nr.} & \textbf{Die SuS können...} & \textbf{AFB} & \textbf{Diff.} \\
\hline
FZ1a & fünf Fragen mit \textit{vuestro/a/os/as} korrekt formulieren & I & alle \\
\hline
FZ2a & einen Partnerdialog über Madrid führen (mit Partnerkarte S. 136) & II & alle \\
\hline
FZ3a & \textit{ser, estar, hay} kontextabhängig korrekt verwenden & II & alle \\
\hline
FZ4a & einen Blogeintrag über das eigene Stadtviertel verfassen (mind. 5 Sätze) & II--III & alle \\
\hline
\end{tabular}

\vspace{1em}
\textbf{Gruppe B:}

\begin{tabular}{|c|p{7cm}|c|c|}
\hline
\textbf{Nr.} & \textbf{Die SuS können...} & \textbf{AFB} & \textbf{Diff.} \\
\hline
FZ1b & das spanische Schulsystem mündlich erklären & II & alle \\
\hline
FZ2b & das deutsche Schulsystem skizzieren und mit dem spanischen vergleichen & II--III & alle \\
\hline
FZ3b & kausale Konjunktionen (\textit{ya que, porque, mientras}) korrekt verwenden & II & alle \\
\hline
FZ4b & einen FP-Flyer sinngemäß auf Spanisch erklären (Mediación) & III & alle \\
\hline
\end{tabular}

% ============================================================================
% 5. METHODISCHE ANALYSE
% ============================================================================
\section{Methodische Analyse}

\subsection{Phasierung (Doppelstunde = 2 $\times$ 45 Min.)}

\textbf{Grundprinzip:} Parallele Arbeit beider Gruppen mit wechselnder Lehrerbetreuung.

\begin{enumerate}
    \item \textbf{1. Stunde (45 Min):}
    \begin{itemize}
        \item Gemeinsamer Einstieg: Begrüßung, Tagesübersicht, Gruppeneinteilung
        \item Getrennte Arbeit: A beginnt Aufg. 9--10, B beginnt Aufg. 6a
        \item Lehrkraft wechselt alle 10--15 Minuten zwischen Gruppen
    \end{itemize}
    \item \textbf{2. Stunde (45 Min):}
    \begin{itemize}
        \item A: Abschluss Dialog, Blogeintrag (Aufg. 11)
        \item B: Vergleich DE/ES (Aufg. 6b), Mediación (Aufg. 7)
        \item Ergebnissicherung in beiden Gruppen
    \end{itemize}
\end{enumerate}

\subsection{Sozialformen}
\begin{itemize}
    \item \textbf{Plenum:} Nur für Einstieg und Schluss (Effizienz bei heterogener Gruppe)
    \item \textbf{Gruppenarbeit:} A und B arbeiten getrennt, innerhalb peer-gestützt
    \item \textbf{Partnerarbeit:} Aufg. 10 (A): Dialog mit verteilten Karten
    \item \textbf{Einzelarbeit:} Aufg. 9 (A): Schriftliche Fragen; Aufg. 11 (A): Blogeintrag
\end{itemize}

\subsection{Medien und Materialien}
\begin{itemize}
    \item Lehrwerk: A\_tope.com (Cornelsen)
    \begin{itemize}
        \item Gruppe A: S. 30--31 + Partnerkarte S. 136
        \item Gruppe B: S. 102--103
    \end{itemize}
    \item Tafel: Kleine Tafel hinten für Gruppe B (Kreide)
    \item iPad/TV: Backup für Projektion (primär analog)
    \item Arbeitsblätter: Keine zusätzlichen (Buchaufgaben ausreichend)
\end{itemize}

\subsection{Differenzierungsmaßnahmen}

\begin{tabular}{|p{3cm}|p{5cm}|p{5cm}|}
\hline
\textbf{Bereich} & \textbf{Gruppe A} & \textbf{Gruppe B} \\
\hline
Komplexität & Gelenkte Produktion mit Vorgaben & Freie Produktion mit Impulsen \\
\hline
Hilfen & Wortschatzbox, Beispielsätze, Partnerkarte & Nur Wörterbuch erlaubt \\
\hline
Tempo & Mehr Zeit für Aufg. 11 (Schreiben) & Zügiges Arbeiten erwartet \\
\hline
Peer-Support & E.H. unterstützt H.P./A.A. & Gruppe arbeitet selbstständig \\
\hline
\end{tabular}

% ============================================================================
% 6. VERLAUFSPLANUNG
% ============================================================================
\section{Verlaufsplanung}

\textbf{1. Stunde (45 Min.)}

\begin{longtable}{|p{0.8cm}|p{1.3cm}|p{4.5cm}|p{3cm}|p{1.5cm}|p{2cm}|}
\hline
\textbf{Zeit} & \textbf{Phase} & \textbf{Lehrerhandeln} & \textbf{Schülerhandeln} & \textbf{SF} & \textbf{Medien} \\
\hline
\endhead

0--5 & Einstieg & Begrüßung auf Spanisch, Tagesübersicht an Tafel, Gruppen\-einteilung bestätigen & Zuhören, Fragen klären & Plenum & Tafel \\
\hline

5--7 & Übergang & A: Aufg. 9 erklären (vuestro-Fragen); B: Aufg. 6a erklären (Schulsystem) & Buch öffnen, Material bereitlegen & Plenum & Buch \\
\hline

7--17 & Erarbei\-tung A & \textbf{Bei Gruppe A:} Aufg. 9 begleiten, vuestro-Formen an Tafel, Beispiel gemeinsam & 5 Fragen formulieren (schriftlich) & EA & Buch, Heft \\
\hline

17--27 & Erarbei\-tung B & \textbf{Zu Gruppe B:} Aufg. 6a besprechen, Struktur Schulsystem klären, Schlüsselbegriffe & Schulsystem erklären (mündlich), Notizen & GA & Buch, Tafel (hinten) \\
\hline

27--37 & Übung A & \textbf{Zurück zu A:} Aufg. 10 einführen, Partnerkarte S. 136 erklären, Paare bilden & Partnerdialog Madrid üben & PA & Buch, Partnerkarte \\
\hline

37--45 & Übung B & \textbf{Zu Gruppe B:} Aufg. 6b: Deutsches System skizzieren, Vergleichspunkte sammeln & Skizze anfertigen, Vergleich formulieren & EA/GA & Heft, ggf. Tafel \\
\hline
\end{longtable}

\vspace{0.5em}
\textit{Hinweis: Während L bei einer Gruppe ist, arbeitet die andere selbstständig weiter. E.H. kann in Gruppe A Hilfestellung geben.}

\vspace{1em}

\textbf{2. Stunde (45 Min.)}

\begin{longtable}{|p{0.8cm}|p{1.3cm}|p{4.5cm}|p{3cm}|p{1.5cm}|p{2cm}|}
\hline
\textbf{Zeit} & \textbf{Phase} & \textbf{Lehrerhandeln} & \textbf{Schülerhandeln} & \textbf{SF} & \textbf{Medien} \\
\hline
\endhead

0--5 & Übergang & Kurze Pause, Fragen klären, Ausblick 2. Stunde & Aufstehen, kurz bewegen & -- & -- \\
\hline

5--15 & Sicherung A & \textbf{Bei Gruppe A:} 2--3 Dialoge vorstellen lassen, Feedback, ser/estar/hay-Fehler korrigieren & Dialoge präsentieren & Plenum A & -- \\
\hline

15--25 & Sicherung B & \textbf{Zu Gruppe B:} Vergleich DE/ES besprechen, Aufg. 7 (Mediación) einführen & Vergleich vorstellen, Aufg. 7 beginnen & GA & Buch \\
\hline

25--35 & Anwen\-dung A & \textbf{Zurück zu A:} Aufg. 11 (Blogeintrag) anleiten: ``Ya sé'' nutzen, mind. 5 Sätze, ser/estar/hay verwenden & Blogeintrag schreiben & EA & Heft \\
\hline

35--40 & Anwen\-dung B & \textbf{Zu Gruppe B:} Aufg. 7 abschließen, Mediación-Strategien besprechen & FP-Flyer erklären (mündlich oder schriftlich) & EA/PA & Buch \\
\hline

40--45 & Abschluss & Zusammenfassung: Was haben wir heute gelernt? HA ankündigen & Reflektieren, HA notieren & Plenum & Tafel \\
\hline
\end{longtable}

\textbf{Legende:} EA = Einzelarbeit, PA = Partnerarbeit, GA = Gruppenarbeit, SF = Sozialform

% ============================================================================
% 7. ERWARTETE SCHWIERIGKEITEN
% ============================================================================
\section{Erwartete Schwierigkeiten}

\begin{tabular}{|p{4cm}|p{9cm}|}
\hline
\textbf{Schwierigkeit} & \textbf{Reaktion} \\
\hline
\textbf{Ser/estar/hay-Verwechslung} (Gruppe A) & Visuelle Merkhilfe an Tafel: SER = Identität, ESTAR = Ort/Zustand, HAY = es gibt. Chorisches Wiederholen. \\
\hline
\textbf{Vuestro-Kongruenz} (Gruppe A) & Endungen tabellarisch darstellen: -o/-a/-os/-as. Bezugswort markieren lassen. \\
\hline
\textbf{H.P. hat Grundlagenlücken} & E.H. als Peer-Tutor einsetzen (kurz!), reduzierte Aufgabenmenge (3 statt 5 Fragen). \\
\hline
\textbf{A.S. unsicher beim Sprechen} & Ermutigung, zuerst schriftlich vorbereiten lassen, dann mündlich. Positive Verstärkung. \\
\hline
\textbf{A.A. schaut zu E.H.} & Für Aufg. 11 trennen, eigenen Stadtviertel-Text verlangen. Lob für eigene Ideen. \\
\hline
\textbf{Gruppe B arbeitet zu schnell} & Zusatzaufgabe: Vergleich erweitern (weitere Unterschiede finden) oder Aufg. 7 vertiefen. \\
\hline
\textbf{Zeitproblem} & Puffer: Blogeintrag (Aufg. 11) kann als HA fortgesetzt werden. Mediación (Aufg. 7) schriftlich = HA. \\
\hline
\textbf{Lehrer-Wechsel zwischen Gruppen} & Timer nutzen (alle 10--15 Min). Gruppe A: E.H./L.N. als Ansprechpartner. Gruppe B: Selbstständigkeit betonen. \\
\hline
\end{tabular}

% ============================================================================
% 8. HAUSAUFGABE
% ============================================================================
\section{Hausaufgabe}

\subsection{Gruppe A [LJ1]}
\begin{itemize}
    \item Blogeintrag (Aufg. 11) fertigstellen, falls nicht abgeschlossen
    \item Mind. 5 Sätze mit ser/estar/hay
    \item Vokabeln S. 30--31 lernen (Stadtviertel, Orte)
\end{itemize}

\subsection{Gruppe B [LJ2/LJ3]}
\begin{itemize}
    \item Aufg. 7 (Mediación) schriftlich abschließen (ca. 80--100 Wörter)
    \item Konjunktionen \textit{ya que, porque, mientras, además} in Sätzen verwenden
    \item Vokabeln S. 102--103 wiederholen
\end{itemize}

\textbf{Kontrolle:} Stichprobenartig zu Beginn DS 2 (Di 06.01.)

% ============================================================================
% 9. ANLAGEN
% ============================================================================
\section{Anlagen}

\begin{enumerate}
    \item Lehrwerk A\_tope.com: S. 30--31 (Gruppe A), S. 102--103 (Gruppe B)
    \item Partnerkarte S. 136 (für Aufg. 10)
    \item Tafelbild: ser/estar/hay-Übersicht, vuestro-Endungen
    \item Timer für Gruppenwechsel (Smartphone oder Uhr)
\end{enumerate}

\textbf{Hinweis:} Da dies eine reine Bucharbeitsstunde ist, werden keine zusätzlichen Arbeitsblätter benötigt. Die Aufgaben des Lehrwerks sind ausreichend für die Doppelstunde.

% ============================================================================
% 10. DIDAKTIK-SCORE
% ============================================================================
\newpage
\section{Didaktik-Score}

\begin{scorebox}
\textbf{Bewertung der didaktischen Qualität nach 10 Dimensionen}\\
\textbf{Ziel-Score: $\geq$ 9.0 (Premium-Qualität)}

\vspace{1em}
\begin{tabular}{|c|l|c|c|p{4.5cm}|}
\hline
\textbf{\#} & \textbf{Dimension} & \textbf{Gew.} & \textbf{Score} & \textbf{Notiz} \\
\hline
1 & Differenzierung & 15\% & 9/10 & Zwei Niveaugruppen, peer-support, aber wenig individuelle Binnendifferenzierung \\
\hline
2 & Taxonomie (AFB) & 10\% & 9/10 & AFB I--III vertreten, Schwerpunkt AFB II (angemessen für Übungsstunde) \\
\hline
3 & Lernziel-Alignment & 10\% & 10/10 & Direkt aus Rahmenplan MV, alle Ziele operationalisiert \\
\hline
4 & Aktivierung & 10\% & 8/10 & Hohe Eigenaktivität, aber Lehrkraft-Wechsel unterbricht Fluss \\
\hline
5 & Scaffolding & 10\% & 9/10 & Gestufte Hilfen (Partnerkarte, Beispiele), E.H. als Peer-Tutor \\
\hline
6 & Feedback-Qualität & 10\% & 8/10 & Mündliches Feedback in Sicherungsphasen, aber zeitlich begrenzt \\
\hline
7 & Sprachliche Klarheit & 10\% & 10/10 & Altersgerecht, klare Arbeitsaufträge \\
\hline
8 & Motivation \& Relevanz & 10\% & 8/10 & Alltagsbezug (Stadtviertel), aber Bucharbeit weniger motivierend \\
\hline
9 & Inklusion (UDL) & 10\% & 8/10 & LRS-Berücksichtigung, aber wenig visuelle Alternativen \\
\hline
10 & Praktikabilität & 5\% & 9/10 & Realistischer Zeitplan, aber Wechsel zwischen Gruppen anspruchsvoll \\
\hline
\multicolumn{3}{|r|}{\textbf{GESAMT:}} & \textbf{8.8/10} & \textcolor{warningorange}{Iteration empfohlen} \\
\hline
\end{tabular}

\vspace{1em}
\textbf{Bewertungsschwellen:}
\begin{itemize}
    \item[$\Box$] \textcolor{successgreen}{$\geq$ 9.0: \textbf{FREIGABE} (Premium-Qualität)}
    \item[$\boxtimes$] \textcolor{warningorange}{8.0--8.9: Iteration empfohlen}
    \item[$\Box$] 6.0--7.9: Überarbeitung erforderlich
    \item[$\Box$] $<$ 6.0: Grundlegende Neukonzeption
\end{itemize}

\vspace{1em}
\textbf{TOP-3 VERBESSERUNGEN:}
\begin{enumerate}
    \item \textbf{Aktivierung (8/10):} Einstieg mit kurzem Spiel/Warm-up erweitern (z.B. Vokabel-Blitzrunde), um alle SuS sofort zu aktivieren.
    \item \textbf{Feedback-Qualität (8/10):} Schriftliche Feedback-Bögen für Partnerarbeit (Aufg. 10) einführen, damit SuS gegenseitig Rückmeldung geben.
    \item \textbf{Motivation (8/10):} Blogeintrag (Aufg. 11) mit realem Publikum verbinden: ``Stellt euch vor, ihr schreibt für Austauschschüler aus Madrid.''
\end{enumerate}

\vspace{1em}
\textbf{Einschätzung:}
Der Score von 8.8 liegt knapp unter der Freigabegrenze. Für eine reine Bucharbeitsstunde (Fortführung der Lehrwerksarbeit) ist dies akzeptabel. Die Hauptschwäche liegt in der Natur der parallelen Gruppenarbeit: Die Lehrkraft kann nicht beide Gruppen gleichzeitig optimal betreuen. Dies wird durch Peer-Tutoring (E.H.) und klare Arbeitsaufträge kompensiert. Für die Beobachtungsstunde (DS 10) wird eine höhere Aktivierung angestrebt.

\end{scorebox}

\end{document}
